\section{Vektorer og Matricer}
\begin{minipage}{0.35\textwidth}
    \[
        \mathbb{F}^{3 \times 1} =
            \begin{bmatrix}
                a_{1} \\
                a_{2} \\
                a_{3}
                \end{bmatrix}
            \mid a_{1}, a_{2}, a_{3} \in \mathbb{F}
    \]
\end{minipage}
\hfill
\begin{minipage}{0.55\textwidth}
    \[
        \mathbb{F}^{1 \times 3} =
            \begin{bmatrix}
                a_{1} & a_{2} & a_{3}
            \end{bmatrix}
            \mid a_{1}, a_{2}, a_{3} \in \mathbb{F}
    \]
\end{minipage}


\subsection{Addition og Subtraktion}
Kun defineret hvis begge matricer har samme størrelse.

\[
    \begin{bmatrix}
        a & b \\
        c & d
    \end{bmatrix}
    +
    \begin{bmatrix}
        w & x \\
        y & z
    \end{bmatrix}
    =
    \begin{bmatrix}
        a + w & b + x \\
        c + y & d + z
    \end{bmatrix}
\]

\[
    \begin{bmatrix}
        a & b \\
        c & d
    \end{bmatrix}
    -
    \begin{bmatrix}
        w & x \\
        y & z
    \end{bmatrix}
    =
    \begin{bmatrix}
        a - w & b - x \\
        c - y & d - z
    \end{bmatrix}
\]

\subsection{Skalar-multiplikation}
\begin{minipage}{0.25\textwidth}
    \[
        k \cdot
        \begin{bmatrix}
            a_{1} \\
            a_{2} \\
            a_{3}
        \end{bmatrix}
        =
        \begin{bmatrix}
            k \cdot a_{1} \\
            k \cdot a_{2} \\
            k \cdot a_{3}
        \end{bmatrix}
    \]
\end{minipage}
\hfill
\begin{minipage}{0.55\textwidth}
    \[
        k \cdot
        \begin{bmatrix}
            a & b & c \\
            d & e & f \\
            g & h & i
        \end{bmatrix}
        =
        \begin{bmatrix}
            k a & k b & k c \\
            k d & k e & k f \\
            k g & k h & k i
        \end{bmatrix}
    \]
\end{minipage}


\subsection{Vektor-multiplikation}
Kun defineret hvis antallet af kolonner i matricen er lig med antallet af rækker i vektoren. Resultatet bliver en vektor med samme antal rækker som matricen.

\[
    \begin{bmatrix}
        a & b & c \\
        d & e & f
    \end{bmatrix}
    \cdot
    \begin{bmatrix}
        x \\
        y \\
        z
    \end{bmatrix}
    =
    \begin{bmatrix}
        ax + by + cz \\
        dx + ey + fz
    \end{bmatrix}
\]

\subsection{Matrix-multiplikation}
Kun defineret hvis antallet af søjler i den første matrix er lig med antallet af rækker i den anden matrix. Resultatet bliver en matrix med antal rækker fra den første matrix og antal kolonner fra den anden matrix.

\[
    \begin{bmatrix}
        w & x \\
        y & z
    \end{bmatrix}
    \cdot
    \begin{bmatrix}
        a & b & c \\
        d & e & f
    \end{bmatrix}
    =
    \begin{bmatrix}
        (wa + xd) & (wb + xe) & (wc + xf) \\
        (ya + zd) & (yb + ze) & (yc + zf)
    \end{bmatrix}
\]


\bemærk{$\mathbf{A} \cdot \mathbf{B} \neq \mathbf{B} \cdot \mathbf{A}$}

\subsection{Algebra}
\subsubsection{Vektorer}
\begin{align*}
    (\vec{u} + \vec{v}) + \vec{w} &= \vec{u} + (\vec{v} + \vec{w}) \\
    \vec{u} + \vec{v} &= \vec{v} + \vec{u} \\
    c \cdot (d \cdot \vec{u}) &= (c \cdot d) \cdot \vec{u} \\
    c \cdot (\vec{u} + \vec{v}) &= c \cdot \vec{u} + c \cdot \vec{v} \\
    (c + d) \cdot \vec{u} &= c \cdot \vec{u} + d \cdot \vec{u}
\end{align*}
\kilde{Sætning 7.1.1}

\subsubsection{Matricer}
%For $\mathbf{A} \in \mathbb{F}^{m \times n}$, $\mathbf{B} \in \mathbb{F}^{n \times l}$ og $\mathbf{C} \in \mathbb{F}^{l \times k}$ gælder:


\begin{align*}
    \mathbf{A}_1 + \mathbf{A}_2 &= \mathbf{A}_2 + \mathbf{A}_1 \\
    (\mathbf{A}_1 + \mathbf{A}_2) + \mathbf{A}_3 &= \mathbf{A}_1 + (\mathbf{A}_2 + \mathbf{A}_3) \\
    \mathbf{A} \cdot (\mathbf{B} \cdot \mathbf{C}) &= (\mathbf{A} \cdot \mathbf{B}) \cdot \mathbf{C} \\
    \mathbf{A} \cdot (\mathbf{B}_1 + \mathbf{B}_2) &= \mathbf{A} \cdot \mathbf{B}_1 + \mathbf{A} \cdot \mathbf{B}_2 \\
    (\mathbf{A}_1 + \mathbf{A}_2) \cdot \mathbf{B} &= \mathbf{A}_1 \cdot \mathbf{B} + \mathbf{A}_2 \cdot \mathbf{B} 
\end{align*}
\kilde{Sætning 7.2.1}

\subsection{Transponering}
\begin{minipage}{0.35\textwidth}
    \[
        \begin{bmatrix}
            a & b & c
        \end{bmatrix}
        =
        \begin{bmatrix}
            a \\
            b \\
            c
        \end{bmatrix}^{T}
    \]
\end{minipage}
\hfill
\begin{minipage}{0.55\textwidth}
\[
    \begin{bmatrix}
        a & b & c \\
        d & e & f \\
        g & h & i
    \end{bmatrix}^{\!T}
    =
    \begin{bmatrix}
        a & d & g \\
        b & e & h \\
        c & f & i
    \end{bmatrix}
\]
\end{minipage}

%For $\mathbf{A} \in \mathbb{F}^{m \times n}$ og $\mathbf{B} \in \mathbb{F}^{n \times l}$ gælder:

\begin{itemize}
    \item $(\mathbf{A}^T)^T=\mathbf{A}$
    \item $(\mathbf{A}_1+\mathbf{A}_2)^T=\mathbf{A}_1^T+\mathbf{A}_2^T$
    \item $(\mathbf{A} \cdot \mathbf{B})^T = \mathbf{B}^T \cdot \mathbf{A}^T $ \hfill \bemærk{Rækkefølgen vendes}
\end{itemize}

\kilde{Sætning 7.2.2}

\subsection{Inverse matricer}
\begin{itemize}
    \item Kun defineret for kvadratiske matricer.
    \item $\mathbf{A}$ har en invers $\iff$ $\rho(\mathbf{A}) =$ antallet af rækker i $\mathbf{A}$.
    \kilde{Lemma 7.3.1}
    \item $\mathbf{A} \cdot \mathbf{A}^{-1}=I_n$ \hfill \kilde{Korollar 7.3.4}
    \item $\mathbf{A}$ har en invers $\iff \det(\mathbf{A}) \neq 0$ \hfill \kilde{Korollar 8.2.4}
\end{itemize}



\subsubsection{Invers ved rækkeoperationer}
For en kvadratisk matrix $\mathbf{A}$ kan $\mathbf{A}^{-1}$ findes ved at bruge rækkeoperationer:
\[
\bigl[\,\mathbf{A} \mid I\,\bigr]
\;\xrightarrow{\text{rref}}\;
\bigl[\,I \mid \mathbf{A}^{-1}\,\bigr].
\]
Kun hvis venstre side ender som identitetsmatricen, eksisterer $A^{-1}$.

\[
    \mathbf{A}=
    \begin{bmatrix}
        a & b \\
        c & d
    \end{bmatrix}
    \implies
    \begin{bmatrix}
        a & b & 1 & 0 \\
        c & d & 0 & 1
    \end{bmatrix}
    \xrightarrow{\text{rref}}
    \begin{bmatrix}
        1 & 0 & x & y \\
        0 & 1 & z & w
    \end{bmatrix}
    \implies
    \mathbf{A}^{-1}=
    \begin{bmatrix}
        x & y \\
        z & w
    \end{bmatrix}
\]

\kilde{Sætning 7.3.2}

\subsubsection{Invers ved detminant}
\[
\begin{bmatrix}
    a & b \\
    c & d
\end{bmatrix}
^{-1}
=
\frac{1}{\det(\mathbf{A})} \cdot
\begin{bmatrix}
    d & -b \\
    -c & a
\end{bmatrix}
\kildetag{Sætning 8.1.1}
\]

\subsection{Lineær afhængighed}
Hvis en eller flere vektorer i et sæt kan skrives som en lineær kombination af de andre vektorer i sættet, er vektorerne lineært afhængige. Ellers er de lineært uafhængige.

\subsubsection{Determinant}
Hvis $\det(\mathbf{A}) = 0$ er søjlerne i $\mathbf{A}$ lineært afhængige. Ellers uafhængige.

\kilde{Korollar 8.2.5}

\subsubsection{Samme rang som vektorer}
Hvis matricen skabt af vekotrerne som søjler har samme rang som antallet af vektorer, er de lineært uafhængige. Ellers er de lineært afhængige.

\kilde{Sætning 7.1.3}
